\section{What the Greek and Ancient Versions Do and Do Not Say}

\subsection{\emph{Othonia}: linen burial cloths}

John uses \emph{othonia} at the burial and at the empty tomb. The term can
denote linen cloths or wrappings. It does not specify a single mummy-like
shroud, a modern fitted garment, or the exact number and dimensions of pieces.
John 19:40 says that Joseph and Nicodemus bound the body with linen cloths and
spices according to Jewish burial custom. Any reconstruction of rigid
wrappings, a hollow body-form, or cloths collapsed in place adds more physical
detail than the noun supplies.

\subsection{\emph{Soudarion}: the face-cloth}

The \emph{soudarion} is the smaller cloth associated with the head or face.
The Latin-derived word is related to sweat-cloth; its New Testament uses include
a cloth employed to wrap money and cloths associated with a person's body.
In John 11:44 it is around Lazarus's face. The Authorized Version's historical
English ``napkin'' does not make it a dining napkin or connect it to a table
etiquette code.

\subsection{\emph{Entetyligmenon}: wrapped or rolled}

The participle describing the face-cloth belongs to a verb used elsewhere of
wrapping Jesus's body in linen. Its lexical range supports ``wrapped up'' or
``rolled up.'' English ``folded'' is possible as a result description in some
translations, but the Greek does not specify a deliberate set of parallel
folds or a culturally coded shape.

The verbal link to burial wrapping creates an important alternative to the
image of Jesus tidying the tomb: the face-cloth may be described in the state
produced by having been wrapped around the head. John tells where and how it
lay relative to the other cloths; he does not narrate an act of post-
Resurrection housekeeping.

\subsection{The Vulgate and the Septuagint boundary}

The Clementine Vulgate preserves the same narrative distinctions without
turning them into a diagram. At the burial it says that the body was bound
with \emph{lintea}; at the tomb it calls the cloths \emph{linteamina} and the
head-cloth a \emph{sudarium}. The latter is \emph{separatim involutum in unum
locum}---apart, wrapped in one place. Luke likewise says that Peter saw the
\emph{linteamina} alone and departed \emph{mirans}, wondering. This Latin
witness therefore reinforces the subject sequence and the language of
wrapping, separation, and wonder, but supplies no folded-table signal or
precise cloth geometry.

The Septuagint bears here only as lexical background: this is a New Testament
scene, not an Old Testament variant. Swete's Greek text of Judges 14:12--13
uses \emph{sindonas} for the thirty garments Samson offers and
\emph{othonia} for the thirty garments his hearers must give if they cannot
solve the riddle; the Alexandrinus apparatus instead reads \emph{sindonas} in
verse 13. This direct locus confirms an ordinary garment use and textual
variation within the Greek tradition. It supplies no burial practice, dining
custom, or coded fold.

\subsection{\emph{Keimena}, \emph{ch\=oris}, and \emph{eis hena topon}}

The cloths are ``lying.'' The face-cloth is not lying with them but
\emph{ch\=oris}, apart, and \emph{eis hena topon}, in one place. The language
establishes distinction and location. It does not yield precise geometry.
Ancient and modern interpreters have imagined the face-cloth above, beside, or
in a separate area; the prepositions do not select a diagram.

\subsection{Three verbs of seeing}

John varies his verbs. The other disciple first \emph{blepei}, sees from the
doorway. Peter \emph{the\=orei}, observes or contemplates, as his attention
moves over the evidence. The other disciple then \emph{eiden}, saw, and
believed. Lexical nuance supports the narrative's staged perception, but Greek
aspect cannot by itself disclose the private content of either man's mind.
John supplies that content only for the final subject.

